\documentclass[11pt]{article}
\usepackage[a4paper,margin=1in]{geometry}
\usepackage[T1]{fontenc}
\usepackage{amsmath,amsthm}
\usepackage{newtxtext,newtxmath}
\usepackage{microtype}
\usepackage{booktabs,tabularx}
\usepackage[font=small,labelfont=bf]{caption}
\usepackage{enumitem}
\usepackage{xcolor}
\usepackage[colorlinks=true,linkcolor=black!70!blue,citecolor=black!70!blue,urlcolor=black!70!blue]{hyperref}
\usepackage[capitalise,noabbrev]{cleveref}

\setlist{itemsep=2pt,topsep=4pt}
\newtheorem{theorem}{Theorem}
\theoremstyle{definition}
\newtheorem{definition}{Definition}
\newcommand{\ip}[2]{\langle #1,#2\rangle}
\newcommand{\norm}[1]{\lVert #1\rVert}
\newcommand{\lean}[1]{\texttt{\small #1}}
\newcommand{\ms}{\,\mathrm{m/s}}
\newcolumntype{Y}{>{\raggedright\arraybackslash}X}

\title{\textbf{Certified Pursuit}\\[4pt]
\large Machine-Checked Necessary and Sufficient Conditions for a Rigid-Body Quadrotor\\
to Keep a Fast Fixed-Wing Aircraft in View}
\author{Abdullah Almanei}
\date{Draft, 2 November 2026}

\begin{document}
\maketitle

\begin{abstract}
We give a full formal proof of a visual pursuit problem in the Lean~4 proof assistant.
The pursuer is a quadrotor with a full rigid-body model in discrete time.
The model has four motors with limits, a mixer, the rotational equation of motion and exact
rotations.
It also has limits on tilt, body rate and speed, parasitic drag and a normal force on the rotor
disc.
A camera on a hinge with a rate limit must keep a fixed-wing target in view in each frame.
The target has limits on speed, turn radius, jerk, snap and climb angle.

We prove an explicit sufficient inequality.
If it is true, a causal method with a delay of four frames keeps each permitted target in view.
The method uses only the positions of the target.
We also prove explicit necessary inequalities: a speed limit for each direction, and a power limit
for the rotor-disc model.
If a necessary inequality is false, a straight and level target escapes from all quadrotor paths.
This is true also for a path that the quadrotor plans with full knowledge of the target's future.

For a quadrotor of 1.5\,kg, the sufficient inequality holds to $52\ms$ and the necessary limit is
$57.7\ms$.
The power limit is within $0.4\ms$ of exact steady force balance.
A point mass that weaves can get to the rest of the gap, but a rigid body cannot fly this weave.
Thus, a better proof must use the rotational limits.
\end{abstract}

\section{Introduction}

Engineers usually design multirotor pursuit of fast aircraft in simulation.
They then test the design in flight.

These two steps cannot answer two important questions.
The first question asks if any method can follow the target with a given airframe.
The second question asks if a given method follows each permitted target.
We answer both questions with proofs in a proof assistant~\cite{lean4,mathlib}.
Our model has sufficient detail to give an answer that an engineer can use.

The paper has six contributions.
\begin{enumerate}
\item A rigid-body pursuit model in Lean~4 (\lean{Model.lean}).
The model has translational and rotational equations of motion in discrete time.
It has exact attitude updates, four motors with limits, a plus mixer, a target class, a camera
hinge and causal methods.
\item A lift theorem (\lean{QuadTraj.ofPlan}).
A kinematic plan is a sequence of positions and attitude frames.
If the thrust, torques, tilt, rate and hinge of a plan are in the limits, the plan is a valid
rigid-body path.
Thus, a proof of sufficiency becomes a set of inequalities on the plan.
\item The sufficient theorem \lean{master\_sufficient3}.
An explicit inequality \lean{MasterIneq3}, which changes with the target speed, makes sure that our
method keeps the target in view.
The method uses a filter with a delay of four frames and a position below the target.
It also uses a frame locked to the side direction and a hinge with limits.
\item The necessary theorems \lean{master\_necessary3} and \lean{master\_necessary\_power}.
These theorems are true for all methods, also for methods with full information.
One target that escapes is sufficient for each proof.
\item Sharp bounds that move the bounds by $10\ms$ to $30\ms$.
The first is an exact formula for the change of body rate.
The others are a sharp bound on the thrust and a time-step term that changes with speed.
\item An analysis of how close the necessary bound is to the physics.
\end{enumerate}

\section{Model}

\subsection{Quadrotor}

The time step is $\delta > 0$.
The state of the quadrotor has six parts.
These parts are the position $x$, the velocity $v$ and an orthonormal frame $(b_1, b_2, b_3 = b_1 \times b_2)$.
They are also the body rate $\Omega$, the motor thrusts $f \in [0, f_{\max}]^4$ and the hinge angle $\eta$.
The translational equation is
\begin{equation}
  v_{k+1} = v_k + \delta\Big[\Big(\tfrac{1}{m}\textstyle\sum_i f_i - k_n\norm{v}\ip{v}{b_3}\Big) b_3
  - g e_3 - c\norm{v} v\Big].
\end{equation}
The rotational equation is
\begin{equation}
  \Omega_{k+1} = \Omega_k + \delta J^{-1}\big(\tau(f) - \Omega\times J\Omega\big).
\end{equation}
An exact rotation through the vector $\delta\,\omega(\Omega)$ updates the frame, with the
Rodrigues formula.
The tilt is not more than $\theta_{\max}$, the rate is not more than $\omega_{\max}$ and the speed
is not more than $v_{\mathrm{env}}$.
The hinge angle and the hinge rate also have limits.

\subsection{Target and visibility}

The target has a speed in the range $[v_{\min}, v_{\max}]$.
Its acceleration along the path and normal to the path has limits, with a minimum turn radius $R$.
Its jerk and snap have limits.
These limits apply to the second and the third differences of velocity.
Its climb angle and vertical acceleration also have limits.
The predicate \lean{StartsSteady} sets the position, velocity and acceleration at the start.

The target is visible in a frame when two conditions are true.
Its distance is not more than the detection range $D$, and its direction is in the field of view
of half-angle $\alpha$.
The predicate \lean{KeepsInView} means that the target is visible in all frames.

\subsection{Methods}

A method, \lean{SteadyStrategy}, gives a quadrotor path for each target path.
It must be causal.
The quadrotor state to step $k+1$ and the motor thrusts to step $k$ can depend only on the target
positions to step $k$.
Our necessary theorems are true for all quadrotor paths.
This is a stronger result than a result for one class of methods.

\section{Results}

\begin{theorem}[Sufficiency, \lean{master\_sufficient3}]
Let \lean{MasterIneq3 P Q C $\delta$ F d} be true, and let the filter $F$ have a delay of four frames.
Then the method keeps each target that starts in steady flight in view.
\end{theorem}

The inequality \lean{MasterIneq3} is a finite list of scalar conditions.
Each condition must be true for each target speed $s$ in the speed range.
The conditions are as follows:
\begin{itemize}
\item the vertical thrust is positive, and the ratio of side thrust to total thrust is less than 1;
\item consecutive frames are near each other;
\item the tilt, the body rate and the hinge rate are in their limits;
\item the motor conditions are true for the roll, pitch and yaw torques, at low and at high thrust;
\item the camera error plus the planner position error is not more than $\alpha$.
\end{itemize}

\begin{theorem}[Speed limit for each direction, \lean{master\_necessary3}, $k_n = 0$]
Let the quadrotor keep each target that starts in steady flight in view.
Then $v_{\max} \le \mathrm{dirSpeed}(u)$ for each permitted direction $u$.
Here $\mathrm{dirSpeed}(u)$ is the highest speed along $u$ with the given drag and tilt limit.
\end{theorem}

\begin{theorem}[Power limit, \lean{master\_necessary\_power}, $k_n > 0$]
Let the quadrotor keep each target that starts in steady flight in view.
Then $B\,v_{\max} \le A$ for each line $A - B s$ that is above
\begin{equation}
  \min\!\Big(F_{\max}s, \frac{F_{\max}^2}{4k_n s}\Big) - c s^3 + \frac{\delta\, \Gamma(s)^2}{2}
\end{equation}
on $[0, v_{\mathrm{env}}]$.
Here $\Gamma(s)$ is a bound on the acceleration at speed $s$.
\end{theorem}

\begin{theorem}[Bounds on both sides, \lean{p3\_bracket} and \lean{p3\_bracket\_power}]
The inequality \lean{MasterIneq3} gives the pursuit result.
The pursuit result gives \lean{NecIneq3} for $k_n = 0$, or \lean{NecPower} for $k_n > 0$.
\end{theorem}

\section{Proof structure}

\subsection{Sufficiency}

The proof of sufficiency has eight steps.

The first step is the planner.
A causal filter uses the past target positions with a delay of four frames.
Filter constants multiplied by $\delta^j$ and the jerk limit give its errors in velocity,
acceleration and position.
A separate bound uses the snap to limit the second difference of the planned acceleration
(\lean{PlannerSnap}).

The second step is the speed window.
Each bound for one step uses a speed window $[s_{\mathrm{lo}}(s), s_{\mathrm{hi}}(s)]$.
This window contains the target speed for three steps (\lean{Under.window}).
With this window, the inequality can be true over a speed range of $20\ms$ to $55\ms$.

The third step defines the frame.
The thrust direction is $n = \mathrm{nrm}(\mathrm{thr})$.
Here $\mathrm{thr}$ is the planned acceleration plus $g e_3$ plus the drag.
The frame follows the side direction $\ell$ of the target:
$b_2 = \mathrm{nrm}(-\ell + \ip{\ell}{n} n)$ and $b_1 = b_2 \times n$ (\lean{LateralFrame}).
This frame exists at all pitch angles if $|\ip{\ell}{n}| < 1$.

The fourth step bounds the body rate.
The body rate is $\delta^{-1}\psi(a)$, with $a$ the antisymmetric part of the frame step
(\lean{BodyRate}).
Two results give sharp bounds.

First, a rotation does not move its own axis.
Thus, the change of body rate is exactly the change of the rotation vector in the new frame
(\lean{RateExact}).
Second, take three consecutive frames $b$, $c$ and $e$.
The change of $a$ is the second difference of the frame in frame $c$, plus products of consecutive
steps $\ip{\Delta' c_i}{\Delta' c_j} + \ip{\Delta c_i}{\Delta c_j}$ (\lean{RateSharp}).
A Parseval identity in the earlier frame bounds these products by tilt multiplied by twist
(\lean{step\_prod}, \lean{twist\_le}).
For the reference quadrotor, this result increased the certified speed from $34\ms$ to $64\ms$ for
dives to $5^\circ$ at jerk $5$.

The fifth step bounds the torque.
The torque is $J\,\Delta\Omega/\delta + \Omega \times J\Omega$, with a bound on each axis.
The motor conditions \lean{mix\_lo} and \lean{mix\_hi} then give the mixer condition
\lean{MixerOK}.

The sixth step bounds the total thrust with the rotor disc.
The total thrust is $\norm{\mathrm{thr}} + k_n\norm{v}\ip{v}{n}$.
The lower bound comes from the power of the planned thrust, $\ip{v_P}{\mathrm{thr}}$.
The upper bound comes from the identity
\begin{equation}
  \norm{\mathrm{thr}}^2 = c^2\norm{v}^4 + 2c\norm{v}\ip{v}{a + g e_3} + \norm{a + g e_3}^2 .
\end{equation}
In this identity, the turn acceleration has no effect, because it is normal to $v$
(\lean{thr\_norm\_le\_Nhi}).
This bound increased the certified speed in the rotor-disc model from $46\ms$ to $52\ms$.

The seventh step is about the hinge and the visibility.
The hinge uses an arctangent law with limits.
The camera error is not more than the bank angle plus the hinge error plus the planner error
(\lean{Hinge}, \lean{UnderTraj.traj\_keeps}).

The eighth step proves causality.
If two target paths agree to step $k$, the plans, frames, rates and motor thrusts also agree
(\lean{Under.traj\_causal}).

\subsection{Necessity}

The proofs of necessity use one target that escapes.
A straight target at $v_{\max}$ escapes from each quadrotor whose displacement along $u$ is not more
than $k\delta V + K$, with $V < v_{\max}$ (\lean{straightTarget\_escapes\_of\_displacement}).

For the speed limit in each direction, the excess of $\ip{v}{u}$ over $V$ decreases by the factor
$1 - 2\delta c V$ at each step (\lean{dir\_excess\_step}).

For the power limit, let $y = \ip{v}{b_3}$.
At each step, the power into the path satisfies
\begin{equation}
  F y - k_n\norm{v} y^2 \le \min\!\Big(F_{\max}\norm v, \frac{F_{\max}^2}{4k_n\norm v}\Big).
\end{equation}
The mechanical energy $\tfrac12\norm{v}^2 + g z$ thus changes by not more than
$\delta(\mathrm{availPower} - c\norm v^3) + \delta^2 \Gamma(\norm v)^2/2$ at each step.
A quadrotor that keeps a level target in view has a height change of not more than $2D$.
Thus, the line condition gives a bound on its displacement along $u$ (\lean{PowerWall}).

\section{How close the bounds are}

\Cref{tab:tight} gives the two bounds for three configurations.

\begin{table}[h]
\centering
\caption{Necessary limit and certified speed at jerk $5$, in m/s.}
\label{tab:tight}
\begin{tabular}{@{}lcc@{}}
\toprule
Configuration & Necessary & Sufficient \\
\midrule
Isotropic drag, dives to $5^\circ$ & $84.5$ & $69.5$ \\
$k_n = 0.01$, dives to $5^\circ$ & $57.7$ & $52.1$ \\
$k_n = 0.01$, dives to $10^\circ$, yaw coefficient $0.03$ & $57.7$ & $52.1$ \\
\bottomrule
\end{tabular}
\end{table}

\subsection{The power limit and force balance}

When the time step goes to zero, the energy bound gives $57.36\ms$ at $k_n = 0.01$.
Exact steady, straight and level flight gives $57.3\ms$.
We also did a numerical search of a Lagrangian dual.
This dual adds multipliers for the mean force balance to the energy term.
The best multiplier was zero.

At $k_n = 0.02$, the bound is $48.2\ms$ and steady flight gives $47.0\ms$.
A point mass can get to this bound with a weave.
It applies more thrust than necessary and turns the extra thrust to the left and to the right in
turn.
The side force does no work, but it tilts the rotor disc away from the airflow.
The point mass gets to $47.8\ms$.
Thus, no bound that uses only translational limits can be lower than approximately $48.2\ms$.

In the full rigid-body model, this weave is not possible.
For the weave, the motors must supply thrusts from $-34$\,N to $+78$\,N, and the limit is 22\,N.
In our search, each weave within the limits was slower than steady flight.
Thus, the rest of the gap comes from the rotational limits.
A better proof must connect three things: the flight direction of the quadrotor, the roll acceleration and
the motor margin.

\subsection{The sufficient side}

The rest of the gap on the sufficient side has two causes.
The first cause is the design of the method.
The method copies the target path exactly, also at short peaks of thrust at top speed.
The second cause is real margin.
The top motor must keep a reserve for attitude control.
Some small terms in the bounds also add to the gap.

\section{Proof method}

The development has approximately 10\,000 lines in 43 files and 522 theorems.
A full build takes approximately five minutes.
A script, \lean{scripts/axcheck.sh}, confirms in continuous integration that each main theorem uses
only the standard axioms.
A blueprint web site links each result to its Lean declarations.

Each inequality \lean{MasterIneq} is a structure of scalar conditions.
A Python file copies each term of this structure (\lean{budget\_v3.py}).
Thus, the numbers and the proofs stay in agreement.
We tested each improvement of a proof with numbers first.
When an improvement did not give a better number, we did not prove it.

We found four useful methods for the proofs:
\begin{enumerate}
\item Give names to terms before you use \lean{linarith}.
This tactic treats products with different syntax as different atoms.
\item Change $k+1+1$ to $k+2$ before you expand plan definitions.
This prevents time-outs when Lean expands definitions.
\item Use explicit chains of \lean{calc} and \lean{mul\_le\_mul} instead of \lean{nlinarith} when
the context is large.
\item Keep abbreviations from \lean{set} in agreement with hypotheses that you add later.
\end{enumerate}

We added the rotor-disc force to the model after the first proofs.
This change had an effect on each user of the translational equation.
The older results now have the explicit hypothesis \lean{Q.normalDrag = 0}.
Thus, their statements did not change without notice.

\section{Limits and future work}

The model does not include the change of motor thrust with airspeed.
It also does not include aerodynamic moments, wind, the ground and sensor noise.
The targets start in steady flight.
A necessary condition for turns is not useful at our ranges.
A quadrotor near the centre of the turn can see a target that flies in a circle.

Two problems are open.
The first problem is a rotational proof for the weave gap.
The second problem is a looser formation that uses range and field of view as margin for agility.
Our results about torque show that physics forces almost no torque in this case.

\section*{Reproducibility}

The repository is \url{https://github.com/AbdullahAlmanei/vision-tracking-lean}.
Use these commands to get the results again:
\begin{itemize}\raggedright
\item Proofs: \texttt{lake build}, then \texttt{scripts/axcheck.sh master\_sufficient3 p3\_bracket p3\_bracket\_power}.
\item Bounds: \texttt{scripts/gap\_v3.py} and \texttt{scripts/gap\_v3.py kn}.
\item Weave study: \texttt{scripts/explore/weave.py}. Power dual: \texttt{scripts/explore/power\_dual.py}.
\end{itemize}

\begin{thebibliography}{9}\small
\bibitem{lean4} L.~de~Moura and S.~Ullrich. The Lean 4 theorem prover and programming language. In
\emph{CADE-28}, 2021.
\bibitem{mathlib} The mathlib Community. The Lean mathematical library. In \emph{CPP}, 2020.
\end{thebibliography}

\end{document}
